\section{Introduction}
\label{sec:introduction}

Federated learning (FL) trains a shared model across decentralised
clients without centralising their data~\cite{mcmahan2017fedavg,
kairouz2021advances}. In the consumer-IoT setting this property is
existentially important: home assistants, smart-home sensors, wearable
biometrics, and connected vehicles continuously emit intimate data
that participants are unwilling, and increasingly not legally
permitted, to upload in raw form. The literature on FL has accordingly
focused on \emph{technical} privacy: how to aggregate updates without
revealing local data, how to limit the per-round information leak, how
to make the global model robust against poisoning.

A different and equally binding privacy axis has been comparatively
under-addressed: \emph{regulatory privacy}. GDPR Article~7 requires
that data-subject consent be demonstrably given, granular over
processing purposes, and freely withdrawable. Article~17 grants the
data subject the right of erasure --- enforcing this on a trained
machine-learning model means the model must \emph{forget} the revoked
participant's contribution. The EU AI Act (Regulation
2024/1689)~\cite{aiact2024} reinforces these obligations for high-risk
systems with Article~15's lifecycle-robustness requirement and
Article~26's deployer-side monitoring obligations. CCPA
\S~1798.105~\cite{ccpa2018} adds the analogous right-to-delete in U.S.
law. None of these regimes assume the FL deployment chooses to comply;
they require demonstrable, audit-grade evidence of compliance.

Three independent research threads have each addressed one corner of
this problem. Consent-management systems exemplified by MedRec
\cite{azaria2016medrec} record consent decisions on a blockchain but
provide no mechanism to influence downstream model training. Machine
unlearning methods such as SISA~\cite{bourtoule2021sisa} and FedEraser
\cite{liu2021federaser} remove a participant's contribution from a
trained model but leave the consent state and the verification step
entirely off-chain. Blockchain-based FL systems
\cite{kim2020blockchained,li2020bflc} log model updates on-chain but
do not enforce consent semantics or verify that unlearning was
actually performed. No published work binds the three threads into a
single chain-enforced protocol --- which is the gap this paper
addresses.

\smallskip
\noindent\textbf{The consent--unlearning--verification triangle.}
We formalise the gap as a missing edge between three established
requirements. Consent (granted, modifiable, revocable, time-bounded,
purpose-limited per GDPR~Art.~7) must remain authoritative over which
data is used in which round. Unlearning (executed by any method ---
SISA, gradient ascent, FedEraser, or future approaches) must be
\emph{verifiable} rather than self-asserted. Verification (against
membership inference, the canonical operational metric for
forgetting~\cite{shokri2017mia,carlini2022lira,choquettechoo2021labelonly})
must in turn drive the consent state, closing the loop. The current
state of the art implements the three vertices but not the edges.

\smallskip
\noindent\textbf{Why on-chain enforcement.} A purely off-chain
consent-and-unlearning workflow trusts the data controller to obey the
data subjects' wishes. Public auditability of consent transitions and
unlearning attestations on a smart-contract substrate removes this
trust. The chain becomes the single source of truth: every consent
state change is a transaction, every unlearning event is gated by a
threshold-signed attestation, every $\dmax$-block timeout produces an
on-chain ComplianceViolation event that downstream monitors can pick
up. The cost of putting this enforcement on-chain is non-trivial
--- a single revocation on Ethereum mainnet at typical gas prices
would cost several dollars. At the $10^3$-device, $10$-operations-per-day
IoT scale this becomes prohibitive. Our experiments show that
deploying the same protocol on a public, gas-free,
EVM-compatible chain with PoA$^2$ consensus reduces this cost to
\emph{exactly} \$0 per operation while preserving every property of
the gas-based deployment.

\smallskip
\noindent\textbf{Contributions.} We make four contributions:
\begin{enumerate}[leftmargin=*]
    \item A \textbf{Consent Lifecycle State Machine} (\clsm) with six
        states, nine transitions, and four formally-stated safety and
        liveness properties. The \clsm is implemented as a soulbound
        ERC-721 contract; the state machine and the on-chain
        representation are isomorphic.
    \item A \textbf{Verifiable Unlearning Protocol} (\vup) that
        accepts an Unlearning Attestation if and only if three
        \emph{independent} membership inference attacks
        (shadow~\cite{shokri2017mia}, LiRA~\cite{carlini2022lira},
        label-only~\cite{choquettechoo2021labelonly}) all score below a
        threshold $\tauMIA$, attested by $t$-of-$n$ ECDSA signatures
        from a registered verifier committee. The protocol is
        method-agnostic over the unlearning algorithm.
    \item A \textbf{three-contract Solidity architecture}
        (ConsentRegistry, UnlearningOracle, FLGatekeeper) deployed on
        \emph{both} a gas-free chain (PureChain, PoA$^2$) and a
        gas-based chain (Polygon Amoy, PoS), producing the first
        side-by-side cost comparison of consent-aware FL under both
        economic models.
    \item A \textbf{rigorous empirical evaluation} on a 9-client
        N-BaIoT~\cite{meidan2018nbaiot} federated autoencoder
        anomaly detector across three approximate unlearning paradigms
        and an exact-retraining baseline. We measure: per-method MIA
        forgetting effectiveness, on-chain gas and USD cost on both
        chains, per-round validation cost as $N$ scales to 500
        participants, and the utility-recovery trajectory after
        revocation.
\end{enumerate}

The remainder of the paper is organised as follows.
Section~\ref{sec:related} positions the work against the three threads.
Section~\ref{sec:system-model} formalises the system model, threat
model, and the four properties we prove.
Section~\ref{sec:framework} details the \clsm, \vup, and the
three-contract architecture.
Section~\ref{sec:evaluation} reports the empirical results.
Section~\ref{sec:discussion} discusses limitations, regulatory
alignment, and the journal-extension trajectory.
Section~\ref{sec:conclusion} concludes.
